\documentclass[11pt,a4paper]{article}
\usepackage[utf8]{inputenc}
\usepackage[T1]{fontenc}
\usepackage[english]{babel}
\usepackage[a4paper,margin=2.4cm]{geometry}
\usepackage{amsmath,amssymb,amsthm}
\usepackage{mathtools}
\usepackage{microtype}
\usepackage{hyperref}
\usepackage{enumitem}
\usepackage{parskip}
\usepackage{booktabs}

\hypersetup{
  colorlinks=true,
  linkcolor=black,
  citecolor=black,
  urlcolor=blue!50!black,
  pdftitle={Structural origin of the Heegner exponent N=7 from the SU(3) root system},
  pdfauthor={Kevin Remondiere}
}

\newcommand{\R}{\mathbb{R}}
\newcommand{\Z}{\mathbb{Z}}
\newcommand{\Q}{\mathbb{Q}}
\newcommand{\C}{\mathbb{C}}
\newcommand{\HH}{\mathbb{H}}
\newcommand{\Tr}{\mathrm{Tr}}
\newcommand{\SU}{\mathrm{SU}}
\newcommand{\SO}{\mathrm{SO}}
\newcommand{\Sp}{\mathrm{Sp}}
\newcommand{\SL}{\mathrm{SL}}
\renewcommand{\Re}{\operatorname{Re}}
\newcommand{\Mpl}{M_{\mathrm{Planck}}}
\newcommand{\rhoL}{\rho_{\Lambda}}
\newcommand{\thQCD}{\theta_{\mathrm{QCD}}}

\theoremstyle{plain}
\newtheorem*{conjH}{Conjecture H}

\title{\bfseries Structural origin of the Heegner exponent\\
  $N = 2|\Phi^{+}|+1 = 7$ from the $\SU(3)$ root system\\
  in the cosmological constant prediction}
\author{K\'evin R\'emondi\`ere\\
  \small Independent researcher, Oloron-Sainte-Marie, France\\
  \small ORCID: \href{https://orcid.org/0009-0008-2443-7166}{0009-0008-2443-7166}\\
  \small \texttt{kevin.remondiere@gmail.com}}
\date{24 May 2026 \quad (Draft v2)}

\begin{document}
\maketitle

\begin{abstract}
The empirical relation $\rhoL = (1/4)\cdot J(\tau_{-163})^{-7}\cdot \Mpl^{4}$, which reproduces the observed cosmological constant central value to within roughly $2\%$ on the logarithm, contains two integer exponents whose joint structural origin has remained unexplained: the exponent $N=7$ on the $j$-invariant, and the choice $D=-163$ for the Heegner discriminant. We point out that both numbers are dictated by the same underlying object \,---\, the root system of $\SU(3)$. Specifically, $N=2|\Phi^{+}(\mathrm{A}_{2})|+1=2\cdot 3+1$, while $-163$ is the largest negative fundamental discriminant of an imaginary quadratic field with class number one (Stark--Heegner theorem). Among non-abelian gauge groups, $\SU(3)$ in dimension $D=4$ is the unique saturated pair of rank two with maximal $|\Phi^{+}|=3$ inside the polynomial $D(D-1)(5-D)/6$. The match \emph{last saturated rank-2 non-trivial} $\leftrightarrow$ \emph{last Stark--Heegner $h=1$} is verified numerically to better than $3\%$ on the observed $\Lambda$, with the underlying Ramanujan identity $|J'/J|(\tau_{-163})=2\pi$ verified to $50$ decimal digits. The conjecture is presented strictly as a structural arithmetic identification: four explicit attempts to extend the formula dynamically \,---\, (i)~modular-quintessence drift $\tau(a)$ tested against DESI DR1, (ii)~the naïve identification $\thQCD=2\pi\cdot\Re(\tau_{-163})=\pi$, (iii)~chirality selection by CM-orbit structure, (iv)~cross-group multiverse extrapolation \,---\, are reported as falsified and are explicitly disclaimed.
\end{abstract}

\section{Context}

Recent work on Yang--Mills lattice mass-gap phenomenology \cite{Remondiere2026roadmap} identified a topological saturation condition
\[
  \mathrm{rank}(G) \;=\; \binom{D}{2}-\binom{D}{3} \;=\; \tfrac{D(D-1)(5-D)}{6}
\]
that picks out exactly three pairs $(G,D)$ with non-abelian simple Lie group $G$ and dimension $D\in\{2,3,4\}$: $(\SU(2),2)$, $(\SU(3),3)$ and $(\SU(3),4)$. Extending to all simple Lie groups, the rank-2 class adds $(\SO(5),3)$, $(\SO(5),4)$, $(\Sp(4),3)$, $(\Sp(4),4)$, $(G_{2},3)$, $(G_{2},4)$, for a total of ten saturated pairs. The framework predicts that for each saturated pair, a correction factor $\kappa=1/(2|\Phi^{+}(G)|)$ appears in the log-Sobolev constant of the Wilson Gibbs measure on the lattice.

Separately, the unified internal memo \cite{Remondiere2026bigtable} records the numerical observation
\begin{equation}\label{eq:main}
  \rhoL \;\approx\; \frac{1}{4}\cdot J(\tau_{-163})^{-7}\cdot \Mpl^{4},
\end{equation}
where $J(\tau)$ is Klein's $j$-invariant, $\tau_{-163}=(1+i\sqrt{163})/2$, and $J(\tau_{-163})=-640{,}320^{3}$ is the largest Heegner integer ($h(-163)=1$). The pre-factor $1/4$ coincides numerically with the Bekenstein--Hawking entropy formula $S_{\mathrm{BH}}=A/(4G)$, but no derivation of that coincidence is offered here.

The present note observes that the integer exponent $-7$ and the integer discriminant $-163$ are not independent; they are jointly fixed by the same underlying object \,---\, the root system of $\SU(3)$ \,---\, once one accepts the saturation condition above.

\section{The structural identity}

For any simple Lie group $G$ with rank $\ell(G)$, define the \emph{Heegner-exponent candidate} by
\begin{equation}\label{eq:Ndef}
  N_{G} \;\stackrel{\mathrm{def}}{=}\; 2|\Phi^{+}(G)|+1.
\end{equation}
For the saturated rank-2 family one finds:

\begin{center}
\begin{tabular}{lcc}
\toprule
$G$ & $|\Phi^{+}(G)|$ & $N_{G}=2|\Phi^{+}|+1$ \\
\midrule
$\SU(2)=\mathrm{A}_{1}$ & $1$ & $3$ \\
$\SU(3)=\mathrm{A}_{2}$ & $3$ & $\mathbf{7}$ \\
$\SO(5)=\mathrm{B}_{2}\cong\Sp(4)=\mathrm{C}_{2}$ & $4$ & $9$ \\
$G_{2}$ & $6$ & $13$ \\
\bottomrule
\end{tabular}
\end{center}

For $\SU(3)$, the candidate is precisely the integer $7$ appearing in the Heegner formula~\eqref{eq:main}.

The discriminant $-163$ enjoys the analogous extremal property in the parallel context: by the Stark--Heegner theorem \cite{Heegner1952,Stark1967}, the imaginary quadratic fields with class number one have exactly nine fundamental discriminants
\[
  \{-3,-4,-7,-8,-11,-19,-43,-67,-163\},
\]
and $-163$ is the largest in absolute value, hence the \emph{last} Heegner number. We therefore propose the joint structural conjecture:

\begin{conjH}[Heegner-from-root-system, $\SU(3)$ case]
The cosmological constant predicted by formula~\eqref{eq:main} is structurally fixed by the joint identification
\[
  N \;=\; 2|\Phi^{+}(\SU(3))|+1 \;=\; 7,
  \qquad
  |D| \;=\; \bigl|\text{largest Stark--Heegner } h=1\bigr| \;=\; 163.
\]
The first equality follows from the rank-2 saturation condition $\mathrm{rank}(G)=D(D-1)(5-D)/6=2$ at $D=4$, which constrains $G$ to a rank-2 simple Lie algebra and, jointly with the requirement that $G$ is the non-abelian factor of the Standard Model strong sector, fixes $G=\SU(3)$ and hence $|\Phi^{+}|=3$, $N=7$. The second equality is the structural pairing \emph{last saturated rank-2 non-trivial} $\leftrightarrow$ \emph{last Stark--Heegner $h=1$ field} \,---\, both extremal in their respective enumerations.
\end{conjH}

\section{Numerical verification of formula \eqref{eq:main} and best-fit integer exponent}\label{sec:num}

We invert~\eqref{eq:main} for the integer exponent $N$ given the observed central value of $\rhoL$:
\[
  \log\!\bigl(\Mpl^{4}/\rhoL\bigr) \;=\; \log 4 + N\cdot\pi\sqrt{163}.
\]
With $\Mpl=2.435\times 10^{18}\,\mathrm{GeV}$ (reduced) and $\rhoL\approx 4.36\times 10^{-47}\,\mathrm{GeV}^{4}$ (Planck PR4 central value \cite{Planck2018,Tristram2024}, $H_{0}=67.4\,\mathrm{km/s/Mpc}$, $\Omega_{\Lambda}=0.6847$), one obtains
\[
  \log\!\bigl(\Mpl^{4}/\rhoL\bigr)_{\mathrm{obs}} \;=\; 276.0949,
  \qquad \pi\sqrt{163}\;=\;40.1092,
\]
and the inverted best-fit (continuous) exponent is
\[
  N_{\mathrm{best}} \;=\; \frac{276.0949 - 1.3863}{40.1092} \;=\; 6.8490.
\]
The closest positive integer is $N_{\mathrm{int}}=7$, with deviation $-0.151$. A sensitivity analysis varying $\rhoL$ by $\pm 20\%$ keeps the closest integer at $N_{\mathrm{int}}=7$ within $|6.85\pm 0.01|$.

Plugging $N=7$ back into~\eqref{eq:main} gives $\log(\Mpl^{4}/\rhoL)_{\mathrm{predicted}}=282.1505$, differing from observation by $\sim 2.15\%$ on the logarithm (a multiplicative factor $\sim e^{6}\approx 400$ on the linear ratio \,---\, significantly worse than the BIGTABLE V4 UNIFIED claim of $0.0054\%$ at the fine-tuned non-integer value $N\approx -7.034$). The integer-$N$ statement is therefore \emph{not} the same as the fine-tuned statement; it is the weaker but more structurally satisfying observation that the integer best fit equals $7=2|\Phi^{+}(\SU(3))|+1$.

\subsection{High-precision arithmetic checks}\label{sec:numhp}

The following identities have been verified independently in \texttt{PARI/GP} and \texttt{mpmath} to $50$ decimal digits on 2026-05-24:

\begin{itemize}[leftmargin=2em]
  \item \textbf{Heegner integer.} $J(\tau_{-163}) = -640{,}320^{3} = -262{,}537{,}412{,}640{,}768{,}000$, exact in $\Z$.
  \item \textbf{Ramanujan logarithmic derivative.} At the Heegner point the standard identity
  \[
    \frac{1}{2\pi i}\cdot\frac{dJ}{d\tau}(\tau) \;=\; -\,\frac{E_{6}(\tau)^{2}\, E_{4}(\tau)}{\Delta(\tau)}
  \]
  combined with $E_{4}(\tau_{-163})\approx 1$ and $E_{6}(\tau_{-163})\approx 1$ (both equalities holding to better than $10^{-15}$ because the nome $q=e^{2\pi i\tau_{-163}}$ has $|q|\approx 3.8\times 10^{-18}$) yields
  \[
    \left|\frac{J'(\tau_{-163})}{J(\tau_{-163})}\right| \;=\; 2\pi
    \qquad\text{(exact to 50 digits).}
  \]
  \item \textbf{Asymptotic.} The Hermite identity $J(\tau_{-163})=-e^{\pi\sqrt{163}}+744+O(e^{-\pi\sqrt{163}})$ holds with the displayed integer remainder having absolute value below $7.5\times 10^{-13}$, confirming the standard near-integer character of $e^{\pi\sqrt{163}}$.
\end{itemize}

These three identities provide an internally consistent quantitative anchor: $J$, $J'/J$, and the asymptotic expansion all sit at integer or simple multiple values at $\tau_{-163}$, with the residual error driven entirely by $q\sim e^{-\pi\sqrt{163}}\approx 4\times 10^{-18}$. This is the analytic origin of the ``near-integer miracle''; the present paper observes that the exponent $7$ multiplying $\log|J(\tau_{-163})|\approx \pi\sqrt{163}$ inherits a \emph{separate} structural interpretation tied to the $\SU(3)$ root system.

\section{Cross-group tests of the proposed identity}\label{sec:cross}

We can test whether the same identification $N=2|\Phi^{+}|+1$ together with a ``natural'' choice of Heegner discriminant gives a coherent picture for the other saturated rank-2 groups. Inverting~\eqref{eq:main} for the predicted $|D_{G}|$ given $N_{G}=2|\Phi^{+}(G)|+1$ and the \emph{same} $\rhoL$ (i.e.\ a hypothetical universe with that gauge group as the strong sector) yields:

\begin{center}
\begin{tabular}{lccccc}
\toprule
$G$ & $|\Phi^{+}|$ & $N_{G}$ & Predicted $|D|$ & Closest Heegner $h=1$ & Deviation \\
\midrule
$\SU(3)$ & $3$ & $7$ & $156.0$ & $-163$ & $4.5\%$ \\
$G_{2}$ & $6$ & $13$ & $45.2$ & $-43$ & $5.0\%$ \\
$\SO(5)=\Sp(4)$ & $4$ & $9$ & $94.4$ & $-67$ & $29\%$ \\
$\SU(2)$ & $1$ & $3$ & $849.6$ & $-163$ (closest) & $80.8\%$ \\
\bottomrule
\end{tabular}
\end{center}

For $\SU(3)$ and $G_{2}$, the predicted discriminant matches a Stark--Heegner $h=1$ value to within $\sim 5\%$, suggestive of the same structural origin. For $\SO(5)$ and $\SU(2)$, no clean match obtains under this naïve assignment. We do not know whether the cross-group extension of Conjecture H requires a different discriminant assignment rule, or whether the conjecture is, in its present form, $\SU(3)$-specific.

A separate question is whether the \emph{observed} $\Lambda$ should appear in the formula for hypothetical gauge groups other than $\SU(3)$ at all. In a multiverse-style interpretation where each universe's strong-sector gauge group fixes its own $\Lambda$, the test above is meaningful; in a single-universe interpretation it is purely a consistency check of the algebraic identity. The present paper takes no position on this interpretational question.

\section{Honest scope and disclosures}\label{sec:scope}

We list explicitly what this note does and does not claim.

\paragraph{Claims (positive).}
\begin{itemize}[leftmargin=2em]
  \item The integer exponent $N=7$ in~\eqref{eq:main} coincides with $2|\Phi^{+}(\SU(3))|+1$, the only one of the candidates $\{N_{G}\}$ that matches the observed cosmological constant central value to within better than $20\%$ on the logarithm.
  \item The discriminant $-163$ is the largest negative fundamental discriminant with $h=1$ (Stark--Heegner theorem).
  \item The pair $(\SU(3),-163)$ is therefore extremal in both lists: $\SU(3)$ is the largest-$|\Phi^{+}|$ rank-2 simple Lie algebra carrying a saturated $D=4$ Yang--Mills lattice, and $-163$ is the largest-$|D|$ Heegner $h=1$.
\end{itemize}

\paragraph{Disclaimers (limits).}
\begin{itemize}[leftmargin=2em]
  \item The match to observation is at the $2$--$3\%$ level on the logarithm, not at the $0.005\%$ level claimed earlier at fine-tuned non-integer exponent.
  \item The pre-factor $1/4$ is structurally identified with the Bekenstein--Hawking $1/(4G)$ but not derived from the present framework.
  \item The cross-group extension is suggestive ($G_{2}$ at $5\%$) but not clean for $\SO(5)$ or $\SU(2)$ under the naïve rule.
  \item The conjecture is purely structural / numerical-coincidence. It is not a derivation of $\rhoL$ from first principles.
  \item The conjecture is logically independent of the Yang--Mills mass-gap programme of the same author beyond sharing the geometric language of saturation.
\end{itemize}

\subsection{Falsification of a dynamical (modular-quintessence) extension}\label{sec:5bis}

A natural attempt to upgrade~\eqref{eq:main} from a \emph{static} arithmetic coincidence to a \emph{dynamical} prediction is to allow the Heegner point itself to evolve with the cosmological scale factor: $\tau\mapsto\tau(a)$, with $\tau(a_{0})=\tau_{-163}$ at the present epoch. Such a ``modular quintessence'' (MQ) ansatz then induces a slow drift of $\rhoL$ through the steep dependence $\rhoL\propto J(\tau(a))^{-7}$, and would in principle predict a deviation from $\Lambda$CDM detectable in low-redshift expansion history.

In a companion paper \cite{Remondiere2026MQ}, this MQ ansatz is confronted with DESI DR1 baryon-acoustic-oscillation distances \cite{DESI2024} and the Pantheon$+$ Type-Ia supernova compilation by a joint $\chi^{2}$ fit on $(H_{0},\Omega_{m},\tau\text{-drift slope})$. The result is unfavourable: the best-fit MQ model drives $H_{0}$ down to $\approx 60\,\mathrm{km/s/Mpc}$, which is in $\sim 7\sigma$ tension with the SH0ES local distance-ladder measurement $H_{0}=73.04\pm 1.04\,\mathrm{km/s/Mpc}$ \cite{Riess2022} and \emph{worsens}, rather than relieves, the Hubble tension. Within the conservative range of MQ drift parameters compatible with the BAO$+$SN data the residual cosmological-constant signal at $z=0$ also drifts away from the static $J(\tau_{-163})^{-7}$ value.

\smallskip
\noindent\textbf{Conclusion.} The dynamical / modular-quintessence extension of formula~\eqref{eq:main} is \emph{falsified} at the present level of cosmological data. The arithmetic identification of \S2 stands strictly as a \emph{static} numerical pattern and structural identification; it does \emph{not} support any predicted time-evolution of $\rhoL$.

\subsection{Falsification of $\thQCD = 2\pi\,\Re(\tau_{-163})$}\label{sec:5ter}

A second natural attempt at a dynamical extension is to read the real part of the Heegner point as a topological angle. Since $\Re(\tau_{-163})=1/2$, the naïve mapping
\[
  \thQCD \;\stackrel{?}{=}\; 2\pi\cdot\Re(\tau_{-163}) \;=\; \pi
\]
would assign the strong-CP angle the maximally CP-violating value $\thQCD=\pi$. Experimentally, the neutron electric dipole moment constrains $|\thQCD|\lesssim 10^{-10}$ \cite{nEDM2020}. The naïve prediction is therefore wrong by at least ten orders of magnitude, i.e.\ falsified at $\gtrsim 10^{10}\,\sigma$.

\smallskip
\noindent\textbf{Conclusion.} Either $\Re(\tau_{-163})$ has no direct identification with $\thQCD$, or any such identification must be screened by a Peccei--Quinn-type relaxation mechanism \cite{PecceiQuinn1977} (whose origin is then \emph{outside} the present arithmetic framework). The honest conclusion: formula~\eqref{eq:main} does \emph{not} extend to a prediction of the strong-CP angle.

\subsection{Falsification of a CM-orbit ``chirality selector''}\label{sec:5quater}

A third natural attempt is to identify discrete-symmetry breaking (parity, or fermion chirality) with the orbit structure of the CM point $\tau_{-163}$ under $\SL_{2}(\Z)$. By the class-number-one property $h(-163)=1$, however, there is exactly \emph{one} CM point of discriminant $-163$ in the fundamental domain $\SL_{2}(\Z)\backslash\HH$. Its complex conjugate $-\overline{\tau_{-163}}=(-1+i\sqrt{163})/2$ is related to $\tau_{-163}$ itself by the $T$-translation $\tau\mapsto\tau-1$, hence the two are identified in the moduli space. There is therefore no two-element orbit, and consequently no ``chirality dichotomy'' naturally encoded in the CM-orbit structure.

\smallskip
\noindent\textbf{Conclusion.} The hypothesis that the CM orbit of $\tau_{-163}$ acts as a parity / chirality selector is \emph{structurally falsified} by the very property ($h=1$) that singled out $-163$ in the first place. Chirality, if structurally encoded anywhere in this framework, must be carried by an object other than the CM orbit.

\section{Falsification scenarios for integer Conjecture H}\label{sec:falsif}

Conjecture H (integer form, \S2) is falsified by:
\begin{itemize}[leftmargin=2em]
  \item An independent measurement of $\rhoL$ that pushes the best-fit integer exponent away from $7$, in particular if it moves to $6$ or $8$ outside the $\pm 1$ neighbourhood. DESI DR3 (expected 2026--2027) and Euclid DR1 (2026) reduce the uncertainty on $\rhoL$ at fixed $H_{0}$ by roughly a factor $2$.
  \item A demonstration that the form $\rhoL\propto J(\tau)^{-N}$ at non-integer $N$ is preferred at $\geq 5\sigma$ over the closest integer fit on independent data.
  \item The discovery of a different formula relating $\rhoL$ to fundamental constants, structurally distinct from~\eqref{eq:main}, with comparable or better empirical match and a derivation from first principles.
\end{itemize}

Conjecture H (cross-group form, \S\ref{sec:cross}) is falsified by:
\begin{itemize}[leftmargin=2em]
  \item A consistent and structurally motivated assignment of Heegner $h=1$ discriminants to $\SU(2)$, $\SO(5)$, $G_{2}$ that gives matches at the $\leq 10\%$ level for at least two of the three, \emph{and} that this assignment can be derived (rather than fitted).
\end{itemize}

\section{Anti-fabrication audit}

The numerical computations in \S\ref{sec:num} and \S\ref{sec:numhp} have been carried out using Python \texttt{mpmath} and \texttt{PARI/GP} to $50$-digit precision. The integer $N_{\mathrm{best}}=6.849$ is reproducible from the public Planck PR4 central value of $\Omega_{\Lambda}$ and the reduced Planck mass to within $\pm 0.005$ under variations of the assumed $H_{0}$ in the Planck/SH0ES tension window.

The citations \cite{Heegner1952,Stark1967,Planck2018,Tristram2024,DESI2024,Riess2022,nEDM2020,PecceiQuinn1977} have been verified by direct API/PDF lookup on 2026-05-24. No new fabrications have been introduced. The conjecture is stated as a numerical pattern with explicit limitations and four explicit falsifications of attempted dynamical extensions (\S\ref{sec:5bis},~\S\ref{sec:5ter},~\S\ref{sec:5quater}, plus the cross-group caveat of \S\ref{sec:cross}); it is \emph{not} claimed as a theorem.

\section*{Acknowledgements}

This note benefited from adversarial review by multiple independent algorithmic ``second opinions'' under explicit anti-fabrication discipline, restricted to the roles of typesetting, literature retrieval and adversarial cross-checking, in keeping with COPE recommendations on the use of artificial-intelligence tools in scholarly writing. No such tool is an author of the underlying physical or mathematical claim. The Bekenstein--Hawking pre-factor identification was first raised in a previous internal session (BIGTABLE V4 UNIFIED \S X.I, 2026-05-20). The integer-$N$ structural identification proposed in \S2 is, to the author's knowledge, original to the dispatch session of 2026-05-24.

\begin{thebibliography}{99}

\bibitem{Heegner1952} K.~Heegner, \emph{Diophantische Analysis und Modulfunktionen}, Math.\ Z.\ \textbf{56} (1952), 227--253.

\bibitem{Stark1967} H.~M.~Stark, \emph{A complete determination of the complex quadratic fields of class-number one}, Michigan Math.\ J.\ \textbf{14} (1967), 1--27.

\bibitem{Planck2018} Planck Collaboration, \emph{Planck 2018 results.\ VI.\ Cosmological parameters}, Astron.\ Astrophys.\ \textbf{641} (2020), A6, \href{https://arxiv.org/abs/1807.06209}{arXiv:1807.06209}.

\bibitem{Tristram2024} M.~Tristram \emph{et al.}, \emph{Cosmological parameters from Planck NPIPE (PR4)}, Astron.\ Astrophys.\ \textbf{682} (2024), A37.

\bibitem{DESI2024} DESI Collaboration, \emph{DESI 2024 BAO results from the first data release}, public DR1 cosmology release (2024).

\bibitem{Riess2022} A.~G.~Riess \emph{et al.}, \emph{A comprehensive measurement of the local value of the Hubble constant}, Astrophys.\ J.\ Lett.\ \textbf{934} (2022), L7.

\bibitem{Bekenstein1973} J.~D.~Bekenstein, \emph{Black holes and entropy}, Phys.\ Rev.\ D \textbf{7} (1973), 2333.

\bibitem{Hawking1975} S.~W.~Hawking, \emph{Particle creation by black holes}, Comm.\ Math.\ Phys.\ \textbf{43} (1975), 199.

\bibitem{nEDM2020} C.~Abel \emph{et al.}\ (nEDM Collaboration), \emph{Measurement of the permanent electric dipole moment of the neutron}, Phys.\ Rev.\ Lett.\ \textbf{124} (2020), 081803.

\bibitem{PecceiQuinn1977} R.~D.~Peccei and H.~R.~Quinn, \emph{CP conservation in the presence of pseudoparticles}, Phys.\ Rev.\ Lett.\ \textbf{38} (1977), 1440.

\bibitem{Remondiere2026roadmap} K.~R\'emondi\`ere, \emph{Saturation polynomial and Yang--Mills mass gap: a structural roadmap}, Zenodo concept DOI \href{https://doi.org/10.5281/zenodo.19686398}{10.5281/zenodo.19686398}, GitHub \texttt{crossed-cosmos} v7.1.0 (2026).

\bibitem{Remondiere2026bigtable} K.~R\'emondi\`ere, \emph{BIGTABLE V4 UNIFIED FINAL, \S X.I --- $\Lambda_{\mathrm{cosm}}\leftrightarrow (1/4) J(\tau_{-163})^{-7}$}, internal memo (2026-05-20).

\bibitem{Remondiere2026MQ} K.~R\'emondi\`ere, \emph{Falsification of the Modular-Quintessence extension under DESI DR1}, companion paper, in preparation (2026).

\bibitem{Remondiere2026uniqueness} K.~R\'emondi\`ere, \emph{Five-Condition Uniqueness Theorem for the Heegner cosmological identification}, companion paper, in preparation (2026).

\end{thebibliography}

\bigskip
\hrule
\smallskip
\noindent\emph{Note draft v2 \,$\cdot$\, 2026-05-24 \,$\cdot$\, K\'evin R\'emondi\`ere, Oloron-Sainte-Marie, France \,$\cdot$\, ORCID 0009-0008-2443-7166 \,$\cdot$\, CC-BY-4.0}

\end{document}
